\section{Switching leaves and completing the proof}
\label{sec:switching}

Fix the tree supplied by Proposition~\ref{prop:tree-data}.
Write $v\preceq s$ when $v=s$ or $v\prec s$.
For $N\geq 1$, set
\[
  \mathcal T_N=\bigcup_{k=0}^{N}L_k,
  \qquad
  B_N(v)=\{s\in L_N:v\preceq s\}
  \quad(v\in\mathcal T_N).
\]
The set $B_N(v)$ is a consecutive block in the ordering of $L_N$, and
\[
  \sum_{s\in B_N(v)}W_s=W_v.
\]
We first show that some path in each finite truncation has all its
functionals large on one vector.

\begin{proposition}[A finite witness]
\label{prop:finite-witness}
For every $N\geq 1$, there are a leaf $s\in L_N$ and a vector $d\in B_Y$
such that
\[
  f_v(d)\geq \frac12
  \qquad
  (r\ne v\preceq s).
\]
\end{proposition}

\begin{proof}
Fix $N$. All the maps, weights, functionals, and budgets in $\mathcal T_N$
are already fixed. Put
\[
  M_k=\min\{e_v:v\in L_1\cup\cdots\cup L_k\}
  \qquad(1\leq k\leq N).
\]

\emph{Choosing a common terminal point.}
For each $v\in\mathcal T_N\setminus\{r\}$, use the parameter
$\theta^v\in\Theta_v$ from Section~\ref{sec:tree}: it assigns the actual
weight $\alpha_h$ to each earlier retained candidate $h\in I_v$, and zero
to each earlier discarded candidate. The point
\[
  q_v=Z_v(\theta^v)
      =\operatorname*{w-lim}_{j\to\mathcal U}Q_v(b_j;\theta^v)
\]
belongs to $E_v$ by \eqref{eq:tree-detector-set}. By \eqref{eq:tree-conditions},
$f_v(y_v-q_v)\geq 1-e_v$. The definition of this weak ultrafilter limit
therefore gives
\[
  \bigl\{j\in\N:
       f_v\bigl(y_v-Q_v(b_j;\theta^v)\bigr)\geq 1-2e_v\bigr\}
       \in\mathcal U.
\]
There are only finitely many such conditions. Moreover,
\eqref{eq:map-displacement} and the approximate fixed point property of
$(b_j)$ give
\[
  \sum_{u\in\mathcal T_N}\norm{G_u(b_j)-b_j}\longrightarrow 0.
\]
Intersecting the finitely many ultrafilter members above with a suitable
cofinite set, choose one $j$ and write $b=b_j$ so that
\begin{align}
  f_v\bigl(y_v-Q_v(b;\theta^v)\bigr)
     &\geq 1-2e_v
     &&(v\in\mathcal T_N\setminus\{r\}),
     \label{eq:terminal-detection}\\
  D_N:=\sum_{u\in\mathcal T_N}\norm{G_u(b)-b}
     &\leq M_N.
     \label{eq:truncation-displacement}
\end{align}

\emph{The switching experiment.}
Initially assign the output $H_s=y_s$ to each leaf $s\in L_N$.
Switch these outputs one at a time to $H_s=b$, proceeding from right to
left. At each stage calculate every other output by
\eqref{eq:tree-rule}. All outputs belong to $K$.
For $v\in\mathcal T_N\setminus\{r\}$, let $U_v$ and $V_v$ be the root
outputs immediately before and immediately after the whole block
$B_N(v)$ is switched.

We need two estimates for subtrees whose leaves are all in the same
state. If none of the leaves below $h\in L_k$ have switched, then
\begin{equation}
  \norm{H_h-y_h}
     \leq \varepsilon_{k,N}:=\sum_{d=k}^{N-1}\tau_d.
  \label{eq:untouched-error}
\end{equation}
Here and below an empty sum is zero. At a leaf this is an equality.
For an internal node, nonexpansiveness and \eqref{eq:child-average}
give the induction step
\begin{align*}
  \norm{H_h-y_h}
  &\leq
    \norm*{\sum_{\widehat u=h}\alpha_u H_u-x}\\
  &\leq
    \sum_{\widehat u=h}\alpha_u\norm{H_u-y_u}
       +\norm*{\sum_{\widehat u=h}\alpha_u y_u-x}
   \leq \varepsilon_{k+1,N}+\tau_k.
\end{align*}
The tolerance schedule \eqref{eq:tree-tolerance} yields, for $k\geq 1$
and every $v\in L_k$,
\begin{equation}
  \varepsilon_{k,N}
    \leq M_k\sum_{d=k}^{\infty}2^{-d-1}
    =2^{-k}M_k
    \leq e_v.
  \label{eq:uniform-level-error}
\end{equation}
Thus the error in \eqref{eq:untouched-error} is at most $e_v$ for every
unchanged node at depth $k$, regardless of which $v\in L_k$ is used in
the comparison.

If all the leaves below $h$ have switched, put
$\delta_u=\norm{G_u(b)-b}$. At each internal node of that subtree,
\[
  \norm{H_u-b}
    \leq \sum_{\widehat w=u}\alpha_w\norm{H_w-b}+\delta_u.
\]
The error at each leaf is zero. Iteration gives
\begin{equation}
  \norm{H_h-b}
    \leq
      \sum_{\substack{u\in\mathcal T_N\setminus L_N\\h\preceq u}}
           \frac{W_u}{W_h}\delta_u
    \leq D_N
    \leq M_N.
  \label{eq:transmitted-error}
\end{equation}
The coefficients $W_u/W_h$ are products of weights below $h$ and are
at most one.

\emph{Detection of an entire block.}
Fix $v\in L_k$, where $1\leq k\leq N$, and consider the stage immediately
before its block switches. Its output is within $e_v$ of $y_v$ by
\eqref{eq:untouched-error}--\eqref{eq:uniform-level-error}.
Let $p_0=r,p_1,\ldots,p_k=v$ be the path to $v$, and put
$A_i=\norm{H_{p_i}-y_v}$ at this stage. Since every off-path child output
and $y_v$ lie in the diameter-one set $K$, nonexpansiveness gives
\[
  A_i\leq
     \alpha_{p_{i+1}}A_{i+1}
       +(1-\alpha_{p_{i+1}})
       +\norm{G_{p_i}(y_v)-y_v}
       \qquad(0\leq i<k).
\]
The off-path contributions telescope:
\[
 \sum_{i=0}^{k-1}W_{p_i}(1-\alpha_{p_{i+1}})
 =\sum_{i=0}^{k-1}(W_{p_i}-W_{p_{i+1}})=1-W_v.
\]
Unwinding the recurrence and using \eqref{eq:tree-conditions}, we obtain
\begin{align}
  \norm{U_v-y_v}
   &\leq
      W_v e_v+(1-W_v)
       +\sum_{i=0}^{k-1}W_{p_i}\norm{G_{p_i}(y_v)-y_v}
       \notag\\
   &\leq 1-W_v+2e_v.
   \label{eq:before-block}
\end{align}

Immediately after the block has switched, precisely the level-$k$
nodes strictly to the left of $v$ have unchanged leaf blocks. By
\eqref{eq:uniform-level-error}, their outputs are within $e_v$ of their
labels $y_h$. The outputs at all the other level-$k$ nodes are within
$M_N\leq e_v$ of $b$ by \eqref{eq:transmitted-error}.
The ideal level-$k$ assignment consisting of these $y_h$ on the left
and $b$ elsewhere has root output $Q_v(b;\theta^v)$ by
\eqref{eq:prediction-cutoff}. The sensitivity estimate
\eqref{eq:tree-sensitivity} and $\sum_{h\in L_k}W_h=1$ therefore imply
\begin{equation}
  \norm{V_v-Q_v(b;\theta^v)}\leq e_v.
  \label{eq:after-block}
\end{equation}
Together, \eqref{eq:terminal-detection}, \eqref{eq:before-block}, and
\eqref{eq:after-block} give
\begin{align}
  f_v(U_v-V_v)
   &\geq f_v\bigl(y_v-Q_v(b;\theta^v)\bigr)
               -\norm{U_v-y_v}-\norm{V_v-Q_v(b;\theta^v)}
               \notag\\
   &\geq (1-2e_v)-(1-W_v+2e_v)-e_v
    =W_v-5e_v.
  \label{eq:block-detection}
\end{align}

\emph{A common witness along one path.}
For a single leaf $s$, only its output changes during its switch.
By \eqref{eq:tree-sensitivity},
\[
  \norm{U_s-V_s}\leq W_s\norm{y_s-b}\leq W_s.
\]
Since $W_s>0$ and both root outputs belong to $K$, the vector
\[
  d_s=\frac{U_s-V_s}{W_s}
\]
belongs to $B_Y$. The successive root differences over a consecutive
block telescope, so
\[
  \sum_{s\in B_N(v)}W_s d_s=U_v-V_v.
\]
Applying $f_v$ to this identity and using \eqref{eq:block-detection}
yields
\begin{equation}
  \sum_{s\in B_N(v)}W_s\bigl(1-f_v(d_s)\bigr)
     =W_v-f_v(U_v-V_v)
     \leq 5e_v.
  \label{eq:weighted-deficit}
\end{equation}
Each summand is nonnegative, because $\norm{f_v}\leq 1$ and
$\norm{d_s}\leq 1$. Summing over the non-root nodes and interchanging the
two finite sums gives
\begin{equation}
  \sum_{s\in L_N}W_s
       \sum_{r\ne v\preceq s}\bigl(1-f_v(d_s)\bigr)
    \leq 5\sum_{v\in\mathcal T_N\setminus\{r\}}e_v
    \leq \frac{5}{128}.
  \label{eq:total-deficit}
\end{equation}
If a leaf $s$ has $f_v(d_s)\leq 1/2$ for any non-root node $v$ on its
path, its inner sum in \eqref{eq:total-deficit} is at least $1/2$. Consequently,
the total weight of all such leaves is at most $5/64$. Since the leaf
weights sum to one, some leaf has $f_v(d_s)>1/2$ for every non-root
node $v$ on its path. This proves the proposition.
\end{proof}

\begin{proof}[Proof of Theorem~\ref{thm:main}]
Suppose a counterexample exists. The reduction and constructions in
Sections~\ref{sec:prelim}--\ref{sec:tree} give the fixed infinite tree to
which Proposition~\ref{prop:finite-witness} applies.
Call a node $v$ \emph{feasible} if some $d\in B_Y$ satisfies
\[
  f_u(d)\geq\frac12\qquad(r\ne u\preceq v).
\]
The root is feasible, and every ancestor of a feasible node is feasible.
Proposition~\ref{prop:finite-witness} gives a feasible node at every depth.
These nodes therefore form a finitely branching tree with arbitrarily
great depth. Starting from the root, choose a child having feasible
descendants at arbitrarily great depths, and repeat. Such a child always
exists because each node has finitely many children. This produces a
branch $r,v_1,v_2,\ldots$, with $v_k\in L_k$, all of whose nodes are
feasible.

For each $k\geq 1$, the set
\[
  C_k=\{d\in B_Y:f_{v_i}(d)\geq 1/2\text{ for }1\leq i\leq k\}
\]
is nonempty and weakly closed in $B_Y$, and $C_{k+1}\subseteq C_k$.
The normalized differences used in Proposition~\ref{prop:finite-witness}
are controlled in the whole unit ball $B_Y$. Reflexivity supplies weak
compactness of this ball, as established in Section~\ref{sec:prelim}.
It follows that there is one vector
\begin{equation}
  d\in\bigcap_{k=1}^{\infty}C_k,
  \qquad
  f_{v_k}(d)\geq\frac12\quad(k\geq 1).
  \label{eq:common-branch-vector}
\end{equation}

On the other hand, the functionals on every branch converge pointwise
to zero on $Y$. Indeed, for fixed $i,j$ and $k\geq\max\{i,j\}$,
\eqref{eq:tree-conditions} and $\diam K=1$ give
\[
  1-e_{v_k}\leq f_{v_k}(y_{v_k}-z_i)\leq 1,
  \qquad
  1-e_{v_k}\leq f_{v_k}(y_{v_k}-z_j)\leq 1.
\]
Subtracting these inequalities yields
\[
  \abs{f_{v_k}(z_i-z_j)}\leq e_{v_k}\longrightarrow 0,
\]
where the convergence follows from the summability of the budgets.
The linear span of the differences $z_i-z_j$ is dense in
$Y=\overline{\operatorname{span}}(K-K)$. Hence the bound
$\norm{f_{v_k}}\leq 1$ extends this convergence to every $d'\in Y$:
for any vector $a$ in that span,
\[
  \abs{f_{v_k}(d')}\leq\norm{d'-a}+\abs{f_{v_k}(a)},
\]
and one first lets $k$ tend to infinity and then lets $a$ approach $d'$.
In particular $f_{v_k}(d)\to 0$, contradicting
\eqref{eq:common-branch-vector}. The assumed counterexample cannot exist.
\end{proof}
